% ============================================================================
%  chapter-template.tex — Tamplate dyal chapitre
%  Bghiti chapitre jdid, dirlih copier coller w semmih
%  bhal: chapters/chapter-01-introduction.tex
% ============================================================================

\chapter{Titre du Chapitre}
\sloppy  % bach mayb9ach derives w ysde3na b hbox

% ── Fo9ra/Section ─────────────────────────────────────────────────────────────────
\section{Première Section}

Rédigez le contenu de votre première section ici. Le template applique
automatiquement la mise en forme définie dans le fichier de style
(\texttt{yazid-template.sty}).

% ── Sous Section ──────────────────────────────────────────────────────────────
\subsection{Sous-section}

Exemple de liste à puces~:

\begin{itemize}

    \item Premier élément de la liste~;

    \item Deuxième élément avec du texte en \textbf{gras} et en
          \textit{italique}~;

    \item Troisième élément.

\end{itemize}

% ── fo9ra khra ─────────────────────────────────────────────────────────
\section{Deuxième Section}

Exemple de liste numérotée~:

\begin{enumerate}

    \item Étape numéro un~;

    \item Étape numéro deux~;

    \item Étape numéro trois.

\end{enumerate}

% ── Mital dyal jdwel ───────────────────────────────────────────────────────────
\section{Exemple de Tableau}

\begin{table}[H]
\centering
\begin{tabularx}{0.9\textwidth}{L{3cm} X C{2.5cm}}
\toprule
\textbf{\color{primarycolor}Colonne A} &
\textbf{\color{primarycolor}Colonne B} &
\textbf{\color{primarycolor}Colonne C} \\
\midrule
Ligne 1 & Description de la ligne 1 & Valeur \\
\addlinespace[0.15cm]
Ligne 2 & Description de la ligne 2 & Valeur \\
\addlinespace[0.15cm]
Ligne 3 & Description de la ligne 3 & Valeur \\
\bottomrule
\end{tabularx}
\caption{Exemple de tableau avec le template}
\end{table}

% ── Mital dyal teswira m9adda ─────────────────────────────────────────────────
\section{Exemple de Figure Encadrée}

\begin{borderedfigure}
    \centering
    % Bdlel had path b lien dyal teswira (f doussi figures/...)
    % \includegraphics[width=0.8\textwidth]{figures/dyalti.png}
    \textcolor{primarycolor}{\rule{0.6\textwidth}{4cm}}  % wa7d lbassa khalouha hna bach tiban test
    \captionof{figure}{Description de la figure encadrée}
\end{borderedfigure}

% ── Mital dyal infobox mzaw9a ────────────────────────────────────────────────────────
\section{Exemple de Boîtes Stylisées}

\begin{infobox}
    \centering
    \vspace{0.3cm}
    {\large Ceci est une \textbf{boîte d'information}. Utilisez-la pour
    mettre en valeur des éléments importants.}
    \vspace{0.3cm}
\end{infobox}

\vspace{0.5cm}

\begin{warningbox}
    \textbf{\textcolor{red!70!black}{\large Attention}}\\[0.3cm]
    Ceci est une boîte d'avertissement pour les informations critiques.
\end{warningbox}

\vspace{0.5cm}

\begin{accentbox}
    \centering
    \vspace{0.3cm}
    {\large Boîte avec bordure dorée pour les formules ou citations.}
    \vspace{0.3cm}
\end{accentbox}
